
\documentclass[a4paper,12pt]{article}

% Sprog og tegnsætning
\usepackage[danish]{babel}
\usepackage[utf8]{inputenc}
\usepackage[T1]{fontenc}

% Layout
\usepackage[a4paper, margin=2.5cm]{geometry}
\usepackage{graphicx}
\usepackage{hyperref}

% Dokumentets metadata
\title{Min rapport}
\author{Dit navn}
\date{\today}

\begin{document}

% Forside
\begin{titlepage}
    \centering

    \vspace*{3cm}

    {\Huge\bfseries Min rapport\par}

    \vspace{1.5cm}

    {\Large Rapportens undertitel\par}

    \vspace{3cm}

    {\large Dit navn\par}

    \vspace{0.5cm}

    {\large Aalborg Universitet\par}

    \vfill

    {\large \today\par}

\end{titlepage}

% Indholdsfortegnelse
\tableofcontents
\newpage

% Kapitel 1
\section{Introduktion}

Dette er introduktionen til rapporten.
Her beskriver du emnet og rapportens formål.

\subsection{Problemstilling}

Her beskriver du problemstillingen.

\subsection{Formål}

Her forklarer du rapportens formål.

% Kapitel 2
\section{Metode}

Her beskriver du de metoder, du anvender.

\subsection{Dataindsamling}

Her beskriver du, hvordan data er indsamlet.

% Kapitel 3
\section{Analyse}

Her skriver du din analyse.

\subsection{Resultater}

Her præsenterer du dine resultater.

% Kapitel 4
\section{Konklusion}

Her samler du op på rapportens vigtigste resultater.

\end{document}
